\section{Introduction}
\label{sec:intro}

Clinical coding assigns standardised diagnostic codes to free-text medical records,
a mandatory step in hospital administration, epidemiological surveillance, and
insurance reimbursement. Manual coding is slow, costly, and prone to inter-annotator
disagreement, motivating the development of automatic approaches \cite{dong2022automated}.

The \textbf{ElCardioCC 2026} shared task targets this problem for Greek cardiology
discharge summaries. Systems must predict all applicable ICD-10 codes from the WHO
10th revision for each document --- a multi-label classification problem over 115 codes.

Three key challenges define the setting. First, \textit{language}: Greek is a
morphologically rich, low-resource language for clinical NLP. Discharge summaries
blend formal Greek medical terminology with transliterated English abbreviations such
as \textit{PCI}, \textit{TAVI}, \textit{CABG}, and \textit{ACS}, creating a
heterogeneous vocabulary that strains standard tokenisers. Second, \textit{document
length}: some summaries exceed the 512-token context limit of standard transformer
encoders, requiring explicit long-document strategies. Third, \textit{label imbalance}:
the 115 ICD-10 codes span frequencies from a single training instance to over 1,000,
with the majority of the label space being rare. A system that ignores rare labels
may still report a high macro-averaged F1 while producing clinically useless
predictions for uncommon diagnoses.

Our approach is driven by the observation that no single model architecture dominates
across all labels and all document types. Frequent labels are best handled by
distributional neural models, while procedural codes with predictable lexical surface
forms are more precisely covered by dictionary-based rules. Rare labels benefit from
high-recall retrieval-based evidence, and mention-level context is critical for
disambiguation between closely related codes (e.g.\ I21 vs.\ I22 vs.\ I25). We
therefore build a system that combines all four paradigms and uses a metaheuristic
ensemble to find the best label-level fusion strategy on the validation set.

The remainder of this paper is structured as follows. Section~\ref{sec:background}
introduces key background concepts. Section~\ref{sec:task} describes the shared task
and dataset. Section~\ref{sec:method} presents the full system methodology.
Section~\ref{sec:results} reports experimental results. Section~\ref{sec:conclusion}
concludes.
